 \begin{center}
        {\large \textbf{Вариант №22 \\
        }}
    \end{center}

\textbf{\underline{Цель работы:} Изучить концепцию виртуальной памяти в Linux, механизмы отображения
виртуальных адресов в физические, организацию памяти процесса, работу подкачки и
кэширования.} 

\begin{enumerate}
    \item \hyperref[sec:1]{Изучить концепцию виртуальной памяти в ОС Linux: назначение, преимущества, отличие от физической памяти.}
    
    \item \hyperref[sec:2]{Ознакомиться с механизмами отображения виртуальных адресов в физические.}
    
    \item \hyperref[sec:3]{Рассмотреть организацию адресного пространства процесса: сегменты кода, данных, кучи, стека, область динамически подгружаемых библиотек (shared libraries).}
    
    \item \hyperref[sec:4]{Освоить использование системных вызовов для управления памятью:}
    \begin{enumerate}[label=\arabic{enumi}.\arabic*]
        \item \hyperref[sec:4.1]{\texttt{mmap()}, \texttt{munmap()} — отображение файлов и анонимное отображение в память.}
        \item \hyperref[sec:4.2]{\texttt{mlock()}, \texttt{munlock()}, \texttt{mlockall()}, \texttt{munlockall()} — блокировка страниц в памяти, предотвращение выгрузки в swap.}
        \item \hyperref[sec:4.3]{\texttt{mprotect()} — изменение прав доступа к страницам (чтение/запись/исполнение).}
        \item \hyperref[sec:4.4]{\texttt{madvise()} — передача советов ядру по предполагаемым шаблонам доступа (например, \texttt{MADV\_SEQUENTIAL}, \texttt{MADV\_RANDOM}, \texttt{MADV\_DONTNEED}); \texttt{posix\_madvise()} — аналогичный механизм в POSIX.}
    \end{enumerate}
    
    \item \hyperref[sec:5]{Изучить организацию взаимодействия функций стандартной библиотеки для динамического выделения памяти (\texttt{malloc()}, \texttt{calloc()}, \texttt{realloc()}, \texttt{free()}) с системными вызовами (\texttt{brk()}, \texttt{mmap()}).}
    
    \item \hyperref[sec:6]{Ознакомиться с механизмами работы ядра с памятью:}
    \begin{enumerate}[label=\arabic{enumi}.\arabic*]
        \item \hyperref[sec:6.1]{Копирование-при-записи (Copy-on-Write, COW) — экономия памяти при порождении процессов через \texttt{fork()}, создание копий страниц только при изменении.}
        \item \hyperref[sec:6.2]{Механизм подкачки страниц — перенос страниц из ОЗУ во вторичное хранилище (swap) и обратно; причины и последствия подкачки.}
        \item \hyperref[sec:6.3]{Кэш страниц — роль в ускорении операций чтения / записи файлов, стратегии замещения страниц.}
        \item \hyperref[sec:6.4]{Обработка отказов по странице (page fault): причины (отсутствие страницы в памяти, защита доступа) и реакции ядра.}
    \end{enumerate}
    
    \item \hyperref[sec:7]{Освоить консольные инструменты для анализа состояния памяти.}
    \begin{enumerate}[label=\arabic{enumi}.\arabic*]
        \item \hyperref[sec:7.1]{\texttt{free} — просмотр использования оперативной памяти и swap-раздела.}
        \item \hyperref[sec:7.2]{\texttt{vmstat} — мониторинг состояния памяти, подкачки и нагрузки на систему.}
        \item \hyperref[sec:7.3]{\texttt{top} / \texttt{htop} — контроль за использованием памяти процессами.}
        \item \hyperref[sec:7.4]{\texttt{/proc/[PID]/maps} и \texttt{/proc/[PID]/smaps} — исследование структуры адресного пространства процесса, распределение анонимной и файловой памяти.}
        \item \hyperref[sec:7.5]{\texttt{/proc/meminfo} — сводная информация об использовании памяти в системе.}
        \item \hyperref[sec:7.6]{\texttt{pmap} — наглядное отображение карт памяти для процесса.}
    \end{enumerate}
    
    \item \hyperref[sec:8]{Выполнить индивидуальное задание для закрепления знаний на практике, подготовив соответствующую программу на языке С (номер задания соответствует номеру студента по журналу; если этот номер больше, чем максимальное число заданий, тогда вариант задания вычисляется по формуле: номер по журналу \% максимальный номер задания, где \% — остаток от деления; если остаток равен 0, студенту назначается задание с максимальным номером).}
\end{enumerate}